\subsection{Dirichlet 问题解的存在性与唯一性}\label{sec:dirichlet-existence-and-uniqueness}

\subsubsection{Dirichlet 问题---解的类型}\label{subsec:dirichlet-types-of-solutions}

先陈述经典 Dirichlet 问题.

\begin{Definition}
\label{def:classical-dirichlet-solution}
设 $\Omega\subset\mathbb{R}^n$ 为有界区域, 给定 $f\in C_b(\Omega)$. 若函数\[u\in C_b^{2k}(\Omega)\cap C_b^{2k-1}(\overline\Omega)\]在 $\Omega$ 中满足
\begin{equation}\label{eq:classical-dirichlet-divergence-equation}
L(\bm{x},D)u=\sum_{0\leq|\sigma|,|\gamma|\leq k}(-1)^{|\sigma|}D^\sigma(a_{\sigma\gamma}(\bm{x})D^\gamma u)=f
\end{equation}
且在 $\partial\Omega$ 上满足
\begin{equation}\label{eq:classical-dirichlet-homogeneous-boundary-conditions}
D^\alpha u=0\quad\text{for }|\alpha|\leq k-1
\end{equation}
则称其为 Dirichlet 问题的\textbf{经典解}.
\end{Definition}

现代分析最重要的思想之一是: 若要保证一个问题的解存在, 通常在一个 ``更大\textquotesingle\textquotesingle\ 的函数空间中更容易做到. 经典 Dirichlet 问题显然如此. 虽然可能期望解具有问题陈述所暗示的全部光滑性, 但首先必须放宽对解的条件, 才能使用前面三章的方法. 放宽解的条件的第一步, 是用 Sobolev 空间陈述问题.

\begin{Definition}
\label{def:strong-dirichlet-solution}
设 $\Omega\subset\mathbb{R}^n$ 为有界区域, 给定 $f\in L^2(\Omega)$. 若函数\[u\in H^{2k}(\Omega)\cap H_0^k(\Omega)\]在 $\Omega$ 中满足
\begin{equation}\label{eq:strong-dirichlet-divergence-equation}
L(\bm{x},D)u=\sum_{0\leq|\sigma|,|\gamma|\leq k}(-1)^{|\sigma|}D^\sigma(a_{\sigma\gamma}(\bm{x})D^\gamma u)=f
\end{equation}
则称其为 Dirichlet 问题的\textbf{强解}.
\end{Definition}

注意以下几点.\begin{enumerate}\item 不仅放宽了对解 $u$ 的条件, 也放宽了对数据 $f$ 的条件. 一旦放宽对 $u$ 的条件, $L^2(\Omega)$ 显然就是 $f$ 所在的自然空间, 因此这种额外的一般性 (包括不连续外力函数等物理上合理的情形) 是 ``免费\textquotesingle\textquotesingle\ 得到的. (事实上, 如下文所示, 每次放宽解的条件时, 都能相应减弱对 $f$ 的条件.)\item 对经典解, 微分方程 \eqref{eq:classical-dirichlet-divergence-equation} 按逐点意义理解. 对强解, 微分方程 \eqref{eq:strong-dirichlet-divergence-equation} 按等价类理解 (方程左右两端表示在 $L^2(\Omega)$ 范数下序列的同一等价类), 或按 ``几乎处处\textquotesingle\textquotesingle\ 的意义理解 (对于学过测度论的读者).\item 不像经典问题那样显式施加边界条件 \eqref{eq:classical-dirichlet-homogeneous-boundary-conditions}, 而是在新问题中将它们纳入空间 $H_0^k(\Omega)$.\item 综合上述观察可见, 新问题确实是经典问题的推广; 即 Dirichlet 问题的每个经典解也都是强解.\end{enumerate}
